% ============================================================
% MATERIALS AND METHODS SECTION - AiN Format
% ============================================================

\section{Materials and Methods}\label{sec:methods}

We formulate sepsis treatment as a finite-horizon Markov Decision Process (MDP) in the offline reinforcement learning setting, where the goal is to learn an optimal policy from a fixed dataset without further environment interaction. We describe the experimental setup including the gym-sepsis simulation environment, offline and online RL algorithms, LEG interpretability analysis, and evaluation protocol.

% ============================================================
\subsection{MDP Formulation}\label{sec:methods:mdp}

The sepsis treatment MDP is defined by the tuple $\mathcal{M} = (\mathcal{S}, \mathcal{A}, \mathcal{P}, \mathcal{R}, \gamma)$. The state space $\mathcal{S} \subset \mathbb{R}^{46}$ captures patient physiological condition through laboratory values, vital signs, and clinical severity scores. The action space $\mathcal{A}$ contains 24 discrete actions indexed from 0 to 23, representing a $5 \times 5$ grid of IV fluid and vasopressor dosing levels. The gym-sepsis environment implements the mapping:
\begin{equation}
a = \min(5 \times \text{IV\_bin} + \text{VP\_bin}, 23)
\end{equation}
effectively merging the highest dosage combination into action 23. The transition dynamics $\mathcal{P}(s_{t+1} | s_t, a_t)$ are learned from MIMIC-III data via the gym-sepsis simulator.$^{1}$ The reward function uses sparse terminal rewards: $\mathcal{R}(s_T, a_T) = +15$ for survival, $-15$ for death, and $0$ for intermediate steps. We use discount factor $\gamma = 0.99$.

A policy $\pi: \mathcal{S} \to \Delta(\mathcal{A})$ maps states to action distributions. The goal is to find $\pi^* = \arg\max_\pi V^\pi(s)$ where the value function is:
\begin{equation}
V^\pi(s) = \mathbb{E}_{\tau \sim \pi, \mathcal{P}} \left[ \sum_{t=0}^{H-1} \gamma^t \mathcal{R}(s_t, a_t) \,\Big|\, s_0 = s \right]
\end{equation}
The action-value function $Q^\pi(s, a)$ represents expected return when taking action $a$ in state $s$ and following $\pi$ thereafter. The optimal policy is derived via $\pi^*(s) = \arg\max_a Q^*(s, a)$.

In offline RL, the agent learns exclusively from a fixed dataset $\mathcal{D} = \{(s_i, a_i, r_i, s'_i)\}_{i=1}^N$ collected under a behavioral policy, without environment interaction during training.$^{2}$ The central challenge is distributional shift, where the learned policy $\pi$ may select out-of-distribution actions with unreliable Q-value estimates due to extrapolation error.$^{3}$

% ============================================================
\subsection{Environment and Data}\label{sec:methods:env}

\subsubsection{Gym-Sepsis Simulator}

We use Gym-Sepsis,$^{1}$ an RL simulator for ICU sepsis treatment trained on MIMIC-III data (Beth Israel Deaconess Medical Center, Boston, MA, USA).$^{4}$

At each timestep, the state is a 46-dimensional vector spanning laboratory values such as lactate, creatinine, and platelet count, vital signs such as blood pressure, heart rate, and SpO$_2$, demographics including age, gender, and race, clinical severity scores including SOFA, LODS, SIRS, qSOFA, and Elixhauser, and treatment status including mechanical ventilation and blood culture. The SOFA score$^{5}$ ranges from 0 to 24, with higher values indicating greater organ dysfunction.

The action space is discrete with 24 actions indexed from 0 to 23, derived from a $5 \times 5$ grid over IV fluid and vasopressor dosage bins. Episodes span ICU stays with 4-hour timesteps until discharge or death. We use sparse rewards: $r_t = +15$ for survival, $-15$ for death, and $0$ for intermediate steps.

\subsubsection{Offline Training Dataset}

We generated an offline dataset of 2,105 episodes with approximately 20K transitions using a heuristic policy based on clinical guidelines,$^{6,7}$ achieving 94.6\% survival. All episodes were used for training without held-out validation or test sets from the offline data. Performance evaluation was conducted by running trained policies on 500 newly sampled episodes in the gym-sepsis simulator.

% ============================================================
\subsection{Algorithms}\label{sec:methods:algos}

We compare 3 offline RL algorithms representing different learning paradigms: Behavior Cloning for supervised learning, Conservative Q-Learning for offline Q-learning, and Deep Q-Network as online RL adapted for offline evaluation. We also evaluate 3 online RL algorithms with architectural innovations. All algorithms use the same neural network architecture for fair comparison, a 3-layer multilayer perceptron with hidden dimensions of 256, 256, and 128 with ReLU activations.

\subsubsection{Offline RL Algorithms}

\paragraph{Behavior Cloning (BC).}
Behavior Cloning treats offline RL as a supervised learning problem, training a policy to imitate the behavioral policy by minimizing the negative log-likelihood of observed actions.$^{8}$ Given a dataset $\mathcal{D} = \{(s_i, a_i)\}_{i=1}^N$ of state-action pairs, BC learns a policy $\pi_{\theta}(a|s)$ by solving:
\begin{equation}
\theta^* = \arg\min_{\theta} -\frac{1}{N} \sum_{i=1}^N \log \pi_{\theta}(a_i | s_i)
\end{equation}
We use d3rlpy's DiscreteBCConfig (Takuma Seno, GitHub, Japan) with batch size 1,024, learning rate $1 \times 10^{-3}$ with Adam optimizer, training for 50,000 gradient steps over 10 epochs.

\paragraph{Conservative Q-Learning (CQL).}
Conservative Q-Learning$^{9}$ is an offline RL algorithm that learns a conservative Q-function to avoid overestimation on out-of-distribution actions. CQL augments the standard Bellman error with a conservatism penalty:
\begin{equation}
\min_Q \alpha \cdot \mathbb{E}_{s \sim \mathcal{D}} \left[ \log \sum_a \exp Q(s, a) - \mathbb{E}_{a \sim \pi_\beta} Q(s, a) \right] + \frac{1}{2} \mathbb{E}_{(s,a,r,s') \sim \mathcal{D}} \left[ (Q(s,a) - \mathcal{T}^\pi Q(s,a))^2 \right]
\end{equation}
where $\alpha$ controls the strength of the conservatism penalty, $\pi_\beta$ is the behavioral policy, and $\mathcal{T}^\pi$ is the Bellman operator. We use d3rlpy's DiscreteCQLConfig with batch size 1,024, learning rate $3 \times 10^{-4}$ with Adam optimizer, $\alpha = 1.0$, target network updates every 2,000 steps, training for 200,000 gradient steps.

\paragraph{Deep Q-Network (DQN).}
Deep Q-Network$^{10}$ combines Q-learning with deep neural networks using experience replay and target networks for stability. The Q-function is updated to minimize the temporal difference error:
\begin{equation}
\mathcal{L}(\theta) = \mathbb{E}_{(s,a,r,s') \sim \mathcal{D}} \left[ \left( Q_\theta(s,a) - \left( r + \gamma \max_{a'} Q_{\theta^-}(s', a') \right) \right)^2 \right]
\end{equation}
We use Stable-Baselines3 (DLR-RM, GitHub, Germany) with batch size 256, learning rate $1 \times 10^{-4}$, target network updates every 1,000 steps, $\epsilon$-greedy exploration from 1.0 to 0.05, training for 100,000 timesteps.

\subsubsection{Online RL Algorithms}

To provide a comprehensive comparison between offline and online RL paradigms, we also evaluate 3 state-of-the-art online RL algorithms with architectural innovations. These algorithms train by interacting with the Gym-Sepsis simulator, collecting 1 million timesteps of experience through exploration.

\paragraph{Double DQN with Attention (DDQN-Attention).}
This algorithm extends Double DQN$^{11}$ with a multi-head self-attention mechanism in the encoder network:
\begin{equation}
h_t = \text{MultiHeadAttention}(s_t, s_t, s_t) + s_t
\end{equation}
The attention mechanism computes scaled dot-product attention across 4 parallel heads, each learning different feature correlations.

\paragraph{Double DQN with Residual Connections (DDQN-Residual).}
This variant incorporates deep residual networks$^{12}$ with skip connections between layers:
\begin{equation}
h_{l+1} = \sigma(\text{LayerNorm}(W_l h_l + b_l + h_l))
\end{equation}
where $\sigma$ is the ReLU activation, $W_l$ and $b_l$ are learnable weights and biases.

\paragraph{Soft Actor-Critic (SAC).}
SAC$^{13}$ is a maximum entropy RL algorithm that optimizes both expected return and policy entropy:
\begin{equation}
J(\pi) = \mathbb{E}_{\tau \sim \pi}\left[\sum_t r(s_t, a_t) + \alpha \mathcal{H}(\pi(\cdot|s_t))\right]
\end{equation}
where $\mathcal{H}(\pi(\cdot|s))$ is the entropy of the policy at state $s$, and $\alpha$ is a temperature parameter. We use the discrete action space variant with a residual encoder architecture.

\paragraph{Training Details.}
All 3 online RL algorithms were trained with 1,000,000 environment interaction steps using experience replay buffers of size 100,000. Training used batch size 256, learning rate $3 \times 10^{-4}$ with Adam optimizer, and target network soft updates with $\tau = 0.005$.

% ============================================================
\subsection{LEG Interpretability Analysis}\label{sec:methods:leg}

To assess interpretability, we employ Linearly Estimated Gradients (LEG),$^{14}$ a model-agnostic perturbation-based method for computing feature importance in RL policies. LEG approximates the policy gradient $\nabla_s Q(s_0, \pi(s_0))$ with respect to state features by locally linearizing the Q-function around a reference state $s_0$.

\subsubsection{LEG Algorithm}

Given a state $s_0 \in \mathbb{R}^{46}$ and a learned policy $\pi$ derived from Q-function $Q(s,a)$ as $\pi(s) = \arg\max_a Q(s,a)$, LEG estimates the saliency vector $\hat{\gamma} \in \mathbb{R}^{46}$ through the following procedure.

For perturbation sampling, we generate $n = 1{,}000$ random perturbations $Z_i \sim \mathcal{N}(0, \sigma^2 I)$ with $\sigma = 0.1$, producing perturbed states:
\begin{equation}
s_i = s_0 + Z_i, \quad i = 1, \ldots, n
\end{equation}

For Q-value evaluation, we compute the Q-value change for each perturbed state $s_i$ and the selected action $a_0 = \pi(s_0)$:
\begin{equation}
\hat{y}_i = Q(s_i, a_0) - Q(s_0, a_0)
\end{equation}

For ridge regression, we compute the sample covariance matrix of perturbations:
\begin{equation}
\Sigma = \frac{1}{n} Z^\top Z + \lambda I, \quad \lambda = 10^{-6}
\end{equation}

For gradient estimation, we estimate the saliency vector via the closed-form ridge regression solution:
\begin{equation}
\hat{\gamma} = \Sigma^{-1} \left( \frac{1}{n} Z^\top \hat{y} \right)
\end{equation}
where $\hat{y} = [\hat{y}_1, \ldots, \hat{y}_n]^\top$. Each element $\hat{\gamma}_j$ quantifies the marginal contribution of feature $j$ to the Q-value.

\subsubsection{State Selection and Analysis Protocol}

We apply LEG to all 6 algorithms using identical parameters for fair comparison. For each algorithm, we analyze 10 representative states sampled from the gym-sepsis environment, ensuring coverage across SOFA severity levels with low being SOFA $\leq$ 5, medium being 6-10, and high being $\geq$ 11.

\subsubsection{Interpretability Metrics}

We quantify interpretability using 3 complementary metrics. Maximum saliency magnitude is the absolute value of the largest saliency score, $\max_j |\hat{\gamma}_j|$, measuring whether any single feature exerts dominant influence on the Q-function. Saliency range is the difference between maximum and minimum saliency scores, indicating how concentrated importance is across features. Clinical coherence is a qualitative assessment of whether top-ranked features align with established medical knowledge, comparing LEG-identified important features against Surviving Sepsis Campaign guidelines.$^{6}$

% ============================================================
\subsection{Evaluation Metrics}\label{sec:methods:eval}

We evaluate algorithm performance using outcome-based and severity-stratified metrics.

\subsubsection{Primary Outcome Metrics}

Survival rate is the proportion of evaluation episodes ending in patient discharge rather than mortality:
\begin{equation}
\text{Survival Rate} = \frac{1}{N} \sum_{i=1}^N \mathbb{1}[R_i > 0]
\end{equation}
where $N$ is the number of evaluation episodes and $R_i = \sum_t r_{i,t}$ is the cumulative return for episode $i$.

Average return is the mean cumulative reward across all episodes:
\begin{equation}
\bar{R} = \frac{1}{N} \sum_{i=1}^N \sum_{t=1}^{T_i} r_{i,t}
\end{equation}

Average episode length is the mean number of timesteps per episode:
\begin{equation}
\bar{T} = \frac{1}{N} \sum_{i=1}^N T_i
\end{equation}

\subsubsection{SOFA-Stratified Analysis}

To assess algorithm performance across patient severity levels, we stratify evaluation episodes by initial Sequential Organ Failure Assessment score.$^{5}$ We define 3 severity strata: Low SOFA ($s_0^{\text{SOFA}} \leq 5$), Medium SOFA ($6 \leq s_0^{\text{SOFA}} \leq 10$), and High SOFA ($s_0^{\text{SOFA}} \geq 11$).

% ============================================================
\subsection{Baseline Policies}\label{sec:methods:baselines}

To contextualize RL algorithm performance, we evaluate 2 baseline policies. The random policy selects actions uniformly at random from the 24-action space at each timestep, with $\pi(a|s) = \frac{1}{24}$. The heuristic policy implements threshold-based rules from sepsis guidelines$^{6}$: it escalates IV fluids when SysBP is below 100 mmHg or lactate is above 2.0 mmol/L, and escalates vasopressors when MeanBP is below 65 mmHg.

% ============================================================
\subsection{Statistical Analysis}\label{sec:methods:stats}

All policies were evaluated on $N = 500$ episodes in the Gym-Sepsis simulator using a standardized protocol. Each episode begins with stochastic initialization sampling an initial patient state from the MIMIC-III-derived distribution. Policies are evaluated deterministically without exploration noise. For discrete policies, we select $a_t = \arg\max_a Q(s_t, a)$. For SAC, we use the mean action.

Episodes terminate when the simulator predicts patient discharge or death, or after 50 timesteps (approximately 8.3 days), whichever occurs first. All results are reported as mean $\pm$ standard deviation. Given the sample size of $N = 500$, performance differences exceeding approximately $\pm 2\%$ for survival rates likely reflect genuine algorithmic differences rather than sampling variability.

All experiments were conducted using Python 3.9 with PyTorch 1.12 (Facebook AI Research, Menlo Park, CA, USA), d3rlpy 1.1.1, and Stable-Baselines3 1.6.0.

% End of Materials and Methods section
